\documentclass[11pt,a4paper]{article}

\usepackage[margin=1in]{geometry}
\usepackage[T1]{fontenc}
\usepackage[utf8]{inputenc}
\usepackage{lmodern}
\usepackage{xcolor}
\usepackage{booktabs}
\usepackage{longtable}
\usepackage{array}
\usepackage{amsmath}
\usepackage{listings}
\usepackage[hidelinks,breaklinks=true]{hyperref}

\definecolor{codebg}{RGB}{248,248,248}
\definecolor{codeframe}{RGB}{220,220,220}
\definecolor{darkblue}{RGB}{0,70,140}
\definecolor{darkgreen}{RGB}{0,110,60}

\lstdefinestyle{cmdstyle}{
  basicstyle=\ttfamily\small,
  breaklines=true,
  columns=fullflexible,
  backgroundcolor=\color{codebg},
  frame=single,
  rulecolor=\color{codeframe},
  keepspaces=true,
  showstringspaces=false
}

\newcommand{\codepath}[1]{\path{#1}}
\newcommand{\pct}[1]{#1\%}
\newcommand{\na}{\textsc{NA}}

\title{Structured Report: Communication and Delay Overhead of Federated Learning with PBC Authentication in LTE VANET}
\author{VANET ns-3 Security Research Workspace}
\date{May 2, 2026}

\begin{document}
\maketitle

\begin{abstract}
This report structures the experimental output comparing two LTE VANET security
cases: a PBC authentication-only baseline and a PBC authentication plus
federated learning (FL) run. The objective is to quantify whether adding FL
introduces extra communication overhead, delay, power consumption, and
security-stage latency. The experiment completed one FL global round with all
selected vehicle updates authenticated successfully. The measured FL-only
communication overhead was very small in this run: 0.006679 J and 2.8768 ms of
analytical communication latency. The current ns-3 FL implementation uses a
lightweight synthetic model-vector training stage, so the measured FL training
compute time is also small; real ML training accuracy and heavier training cost
remain in the Python baseline.
\end{abstract}

\tableofcontents
\newpage

\section{Purpose of the Comparison}

The comparison answers the question:

\begin{quote}
How much extra communication overhead and delay are introduced when federated
learning is enabled on top of the existing PBC-authenticated VANET model?
\end{quote}

Two cases were compared under the same LTE network label:

\begin{longtable}{p{0.28\linewidth}p{0.62\linewidth}}
\toprule
\textbf{Case} & \textbf{Meaning} \\
\midrule
Authentication only & Vehicles use the existing PBC security model for VANET authentication, signing, aggregation, verification, and communication-energy accounting. FL is disabled. \\
Authentication + FL & The same PBC security model is active, and federated learning is enabled. Vehicles receive a model, compute a local synthetic update, mask it, sign it, upload it to RSUs, RSUs verify and edge-aggregate updates, and the BS performs global aggregation. \\
\bottomrule
\end{longtable}

The generated source result files are:

\begin{itemize}
  \item \codepath{simulation_workspace/ns-allinone-3.44/ns-3.44/results/fl_auth_compare/lte_pbc_comparison.md}
  \item \codepath{simulation_workspace/ns-allinone-3.44/ns-3.44/results/fl_auth_compare/lte_pbc_comparison.csv}
  \item \codepath{simulation_workspace/ns-allinone-3.44/ns-3.44/results/fl_auth_compare/lte_pbc_auth_power.csv}
  \item \codepath{simulation_workspace/ns-allinone-3.44/ns-3.44/results/fl_auth_compare/lte_pbc_auth_fl_power.csv}
  \item \codepath{simulation_workspace/ns-allinone-3.44/ns-3.44/results/fl_auth_compare/lte_pbc_auth_stages.csv}
  \item \codepath{simulation_workspace/ns-allinone-3.44/ns-3.44/results/fl_auth_compare/lte_pbc_auth_fl_stages.csv}
\end{itemize}

\section{Experiment Setup}

The test was run from the ns-3 workspace with the LTE and NR security scenarios
already built. The relevant comparison output shows that both LTE runs reached
29 seconds of simulated progress, so the authentication-only and
FL-enabled cases are comparable in simulation duration.

\begin{lstlisting}[style=cmdstyle,caption={Command used to run the LTE FL/authentication overhead comparison.}]
cd /home/sahib/vanet_ns3_security_research
SIM_TIME=30 FL_ROUNDS=3 FL_PARTICIPANTS=10 FL_MODEL_DIM=16 ./scripts/run_fl_auth_overhead_compare.sh lte
\end{lstlisting}

Important parameters used in this comparison are summarized below.

\begin{longtable}{p{0.33\linewidth}p{0.57\linewidth}}
\toprule
\textbf{Parameter} & \textbf{Value / Meaning} \\
\midrule
Network label & \codepath{5g_lte} \\
Security mode & PBC authentication \\
Simulation duration & 30 s requested; progress log reached 29 s before simulator stop \\
FL enabled in second case & Yes \\
FL rounds requested & 3 \\
FL participants & 10 selected vehicles \\
FL model dimension & 16 double-valued model weights \\
FL architecture & Vehicles $\rightarrow$ RSUs $\rightarrow$ BS hierarchical aggregation \\
FL training model in ns-3 & Lightweight synthetic vector-update simulation \\
Real ML baseline & Kept separately in \codepath{Federated-Learning-VANET} \\
\bottomrule
\end{longtable}

\section{Overhead Formula}

For every metric, the report uses:

\[
\text{overhead}_{abs}=\text{value}_{auth+FL}-\text{value}_{auth-only}
\]

\[
\text{overhead}_{\%}=100\times
\frac{\text{value}_{auth+FL}-\text{value}_{auth-only}}
{\text{value}_{auth-only}}
\]

Positive values mean that authentication plus FL is higher than authentication
only. Negative values mean the FL-enabled run measured lower for that metric.
Small negative values can happen because wall-clock cryptographic timing and
packet events have run-to-run variation; they should not be interpreted as FL
magically reducing cryptographic cost.

\section{FL Completion and Authentication Success}

The FL-enabled run completed one full global FL round.

\begin{longtable}{p{0.42\linewidth}p{0.45\linewidth}}
\toprule
\textbf{Metric} & \textbf{Measured value} \\
\midrule
Completed FL rounds & 1 \\
Last completed round index & 0 \\
Completion simulation time & 10.402010 s \\
Selected FL vehicles & 10 \\
RSU count participating in FL & 2 \\
Vehicle updates sent & 10 \\
Vehicle updates verified & 10 \\
Vehicle updates rejected & 0 \\
Authentication success rate for FL updates & $10/10=100\%$ \\
Model download bytes & 1632 bytes \\
Update upload bytes & 1440 bytes \\
Global loss proxy & 0.500000 \\
Global accuracy proxy & 0.500000 \\
\bottomrule
\end{longtable}

\textbf{Interpretation.} All 10 participating FL vehicles successfully submitted
authenticated updates. No FL update was rejected. This means the PBC-secured FL
message path worked for the completed round.

\section{Communication Overhead Results}

\begin{longtable}{p{0.32\linewidth}rrrr}
\toprule
\textbf{Metric} & \textbf{Auth only} & \textbf{Auth + FL} & \textbf{Overhead} & \textbf{Overhead \%} \\
\midrule
Communication energy (J) & 5.538350 & 5.545030 & 0.006680 & 0.120614 \\
Average communication power (W) & 0.190978 & 0.191208 & 0.000230 & 0.120433 \\
Total communication-stage energy (J) & 5.744753 & 5.743032 & -0.001721 & -0.029966 \\
Total communication-stage latency (ms) & 6093.272800 & 6086.996160 & -6.276640 & -0.103009 \\
FL-only communication energy (J) & 0.000000 & 0.006679 & 0.006679 & \na \\
FL-only communication latency (ms) & 0.000000 & 2.876800 & 2.876800 & \na \\
FL communication events & 0.000000 & 92.000000 & 92.000000 & \na \\
\bottomrule
\end{longtable}

\textbf{Interpretation.} The direct periodic communication-energy metric
increased from 5.538350 J to 5.545030 J, which is only 0.006680 J higher. This
is a 0.120614\% increase, so the FL communication overhead is small compared
with the existing VANET communication workload in this LTE run. The FL-specific
communication rows show 92 FL communication events and 2.876800 ms of analytical
airtime latency.

The slight negative value in total communication-stage energy and latency is a
measurement-level difference caused by event counts and run-to-run variation in
the non-FL background traffic. The most direct FL communication overhead is the
FL-only row: 0.006679 J and 2.876800 ms.

\section{Network Delay Results}

\begin{longtable}{p{0.34\linewidth}rrrr}
\toprule
\textbf{Metric} & \textbf{Auth only} & \textbf{Auth + FL} & \textbf{Overhead} & \textbf{Overhead \%} \\
\midrule
Average V2V delay (ms) & 1116.910000 & 1113.690000 & -3.220000 & -0.288295 \\
Average V2I uplink delay (ms) & 19.641300 & 18.927000 & -0.714300 & -3.636725 \\
Average V2I downlink delay (ms) & 1363.880000 & 1368.400000 & 4.520000 & 0.331407 \\
Average V2I RTT (ms) & 1378.000000 & 1382.540000 & 4.540000 & 0.329463 \\
Average registration delay (ms) & 495.800000 & 495.800000 & 0.000000 & 0.000000 \\
\bottomrule
\end{longtable}

\textbf{Interpretation.} The main V2I RTT increased by 4.540 ms, which is only
0.329463\%. V2I downlink delay increased by 4.520 ms, also about 0.33\%. V2V
and uplink delay were slightly lower in the FL run, which is best interpreted
as normal simulation/event variation rather than a true improvement caused by
FL.

The registration delay was unchanged at 495.800 ms. This is expected because FL
starts after credentials are already available, so FL does not change the PBC
registration path in this version.

\section{Security and Power Results}

\begin{longtable}{p{0.34\linewidth}rrrr}
\toprule
\textbf{Metric} & \textbf{Auth only} & \textbf{Auth + FL} & \textbf{Overhead} & \textbf{Overhead \%} \\
\midrule
Total energy (J) & 345.910000 & 312.334000 & -33.576000 & -9.706571 \\
Average total power (W) & 11.927900 & 10.770200 & -1.157700 & -9.705816 \\
Average signing latency (ms) & 25.793100 & 21.298700 & -4.494400 & -17.424815 \\
Average verification latency (ms) & 106.041000 & 96.257900 & -9.783100 & -9.225771 \\
PBC-stage energy (J) & 354.953512 & 318.368395 & -36.585117 & -10.307017 \\
PBC-stage latency (ms) & 24387.827668 & 22199.734695 & -2188.092973 & -8.972070 \\
\bottomrule
\end{longtable}

\textbf{Interpretation.} The FL-enabled run measured lower total security energy
and PBC-stage latency than the authentication-only run. This should be treated
carefully. The PBC signing and verification code is measured using wall-clock
cryptographic timing, so values can vary between runs even when the simulation
configuration is the same. The correct conclusion is not that FL reduces PBC
cost. The correct conclusion is that, for this run, the added FL-specific
computation and communication were small compared with the dominant PBC
aggregate verification workload.

\section{FL-Specific Computation Results}

\begin{longtable}{p{0.34\linewidth}rrrr}
\toprule
\textbf{Metric} & \textbf{Auth only} & \textbf{Auth + FL} & \textbf{Overhead} & \textbf{Overhead \%} \\
\midrule
FL-only computation energy (J) & 0.000000 & 0.000392 & 0.000392 & \na \\
FL-only computation latency (ms) & 0.000000 & 0.100141 & 0.100141 & \na \\
FL computation events & 0.000000 & 42.000000 & 42.000000 & \na \\
\bottomrule
\end{longtable}

The FL-specific computation stages include:

\begin{itemize}
  \item \codepath{fl_local_train}: synthetic local model update at vehicles.
  \item \codepath{fl_mask_generate}: masking step before upload.
  \item \codepath{fl_edge_aggregate}: RSU-side weighted aggregation.
  \item \codepath{fl_global_aggregate}: BS-side global aggregation.
\end{itemize}

\textbf{Interpretation.} In this v1 ns-3 implementation, local FL training is
intentionally lightweight. Therefore, FL-only computation overhead is very small:
0.000392 J and 0.100141 ms. If the project needs to demonstrate heavier FL
training delay, the next step should be to add a configurable simulated training
delay, such as \codepath{--flLocalTrainDelayMs}, or to couple the ns-3 scenario
with measured training profiles from the Python baseline.

\section{Overall Conclusion}

For this LTE PBC comparison, adding FL produced the following key outcomes:

\begin{itemize}
  \item FL authentication succeeded for all selected vehicles: 10 verified out of 10 submitted updates.
  \item FL introduced a small direct communication overhead: 0.006679 J and 2.876800 ms.
  \item Average communication power increased by only 0.120433\%.
  \item V2I RTT increased by 4.540 ms, approximately 0.329463\%.
  \item Registration delay did not change because FL starts after PBC credential setup.
  \item Synthetic FL computation overhead was very small: 0.000392 J and 0.100141 ms.
  \item PBC aggregate verification remains the dominant security cost, not FL aggregation itself.
\end{itemize}

Therefore, in the current implementation, the FL communication overhead is
measurable but small relative to the existing PBC-authenticated VANET workload.
The strongest evidence is the FL-only communication overhead and the
successfully authenticated FL update count. For a more realistic ML-training
latency study, the simulator should either inject a configurable training delay
or import measured training-time profiles from the Python FL baseline.

\section{Recommended Next Experiment}

To show the effect of heavier FL training, run the same comparison with a larger
model dimension and more rounds:

\begin{lstlisting}[style=cmdstyle,caption={Recommended heavier FL comparison command.}]
cd /home/sahib/vanet_ns3_security_research
SIM_TIME=90 FL_ROUNDS=5 FL_PARTICIPANTS=20 FL_MODEL_DIM=64 ./scripts/run_fl_auth_overhead_compare.sh lte
\end{lstlisting}

If training delay should affect simulated network time directly, add a future
scenario option such as:

\begin{lstlisting}[style=cmdstyle,caption={Proposed future option for simulated FL training delay.}]
--flLocalTrainDelayMs=50
\end{lstlisting}

This would make the FL training time visible not only in power/stage CSVs, but
also in the simulated round completion time and end-to-end FL latency.

\end{document}
